\chapter{FUNDAMENTAÇÃO TEÓRICA E REVISÃO BIBLIOGRÁFICA}
\label{chap:fundamentacao_teorica}

O planejamento de Manobras de Evasão de Colisão (\textit{Collision Avoidance Maneuvers} - CAM) representa um problema central na operação de veículos espaciais. Sua resolução integra conceitos de mecânica orbital, modelagem de incertezas e otimização matemática. Este capítulo apresenta a fundamentação teórica e a revisão bibliográfica necessárias para situar o trabalho no estado da arte. A estrutura do texto aborda o contexto do lixo espacial, os fundamentos da dinâmica e análise de risco, e os métodos de otimização de manobras.

\section{Impacto dos Detritos Espaciais nas Operações de Satélites}
\label{sec:ambiente_lixo_espacial}

O ambiente orbital terrestre, notadamente a região de Órbita Baixa (LEO), sofreu alterações substanciais desde o início da era espacial. A região, anteriormente livre de obstáculos artificiais, caracteriza-se hoje por uma alta densidade de objetos, incluindo satélites ativos, cargas úteis desativadas, estágios de foguetes e fragmentos resultantes de colisões e explosões. Esta configuração impõe riscos operacionais contínuos às missões espaciais.

Atualmente, a população de objetos catalogados ultrapassa a ordem de dezenas de milhares. De acordo com dados consolidados pela Agência Espacial Europeia (ESA), as redes de vigilância rastreiam cerca de $40.500$ objetos espaciais. Deste total, aproximadamente $11.100$ correspondem a satélites operacionais, indicando que a maior parte dos objetos em órbita consiste em detritos espaciais. Além dos objetos catalogados, estimativas baseadas em modelos estatísticos apontam para a existência de cerca de $1,1$ milhão de fragmentos com dimensões entre $1$ cm e $10$ cm, não rastreáveis, mas com energia cinética suficiente para comprometer ativos operacionais em caso de impacto \cite{ESA_EnvReport2025}.

Relatórios de agências como a ESA e a NASA documentam o crescimento da população de objetos, impulsionado pelo lançamento de grandes constelações e por eventos de fragmentação \cite{NASA_ODQN_2024}.

Neste contexto, Donald J. Kessler e Burton G. Cour-Palais publicaram, em 1978, um trabalho fundamental descrevendo o fenômeno hoje denominado Síndrome de Kessler \cite{Kessler1978}. Os autores postularam que, ao ultrapassar uma densidade crítica de objetos, a taxa de geração de detritos por colisões em cascata superaria a taxa de decaimento natural por arrasto atmosférico, resultando em um crescimento autossustentável da população de detritos. Estudos subsequentes corroboram essa teoria, sugerindo que determinadas regiões em LEO já operam próximas a este ponto de instabilidade \cite{Liou2010}. A ocorrência de eventos como a colisão entre os satélites Iridium 33 e Cosmos 2251, em 2009, valida concretamente esse risco e fundamenta a necessidade da implementação rotineira de manobras de evasão de colisão nos centros de controle.

A execução de manobras de evasão consolidou-se como um procedimento nominal na operação de satélites. Entretanto, a abordagem tradicional, baseada na análise manual de riscos, apresenta limitações de escalabilidade frente à implementação de grandes constelações. O volume de alertas de conjunção (CDM - \textit{Conjunction Data Messages}) processados diariamente torna proibitiva a manutenção de processos exclusivamente dependentes de intervenção humana, fundamentando a demanda por sistemas de Evasão de Colisão Automatizados (\textit{Automated Collision Avoidance} - ACA) \cite{Virgili2023}.

Tais sistemas devem possuir a capacidade de processar massivamente eventos de conjunção e determinar soluções ótimas. A automação, contudo, introduz desafios específicos: os algoritmos devem operar respeitando restrições operacionais rígidas, como janelas de visibilidade para telecomando, orçamento de propelente e a garantia de que a nova trajetória não resulte em conjunções secundárias \cite{Lemmens2022}. A presente dissertação insere-se neste contexto, propondo métodos de otimização que satisfaçam estes requisitos de segurança e eficiência.

\section{Fundamentos da Mecânica Orbital}
\label{sec:fundamentos_mecanica_orbital}

A determinação precisa da trajetória de satélites fundamenta a análise de risco de colisão. Na dinâmica orbital, a seleção do método de propagação envolve um balanço (\textit{trade-off}) entre a fidelidade do modelo físico e o custo computacional associado. A escolha do propagador deve ser compatível com os requisitos de precisão da missão e com a qualidade dos dados de rastreio disponíveis. Esta seção apresenta os princípios da mecânica orbital, abrangendo desde o problema idealizado de dois corpos até os modelos numéricos de alta fidelidade que incorporam as perturbações ambientais.

\subsection{O Sistema Geocêntrico Inercial (ECI)}
\label{subsec:eci}

Para satélites que orbitam a Terra, o sistema de coordenadas fundamental é o Geocêntrico Inercial (ECI - \textit{Earth-Centered Inertial}). Embora a Terra esteja em movimento ao redor do Sol, para fins de satélites próximos, este referencial é considerado suficientemente inercial (quasi-inercial).

O sistema ECI é definido por uma tríade cartesiana dextrogira ($\hat{I}, \hat{J}, \hat{K}$) com as seguintes características, ilustradas na Figura \ref{fig:eci_frame}:

\begin{itemize}
    \item \textbf{Origem:} Coincide com o centro de massa da Terra.
    \item \textbf{Plano Fundamental (Plano Equatorial):} O plano $XY$ coincide com o plano do equador terrestre.
    \item \textbf{Eixo Principal ($+\hat{K}$):} Aponta para o Polo Norte Celeste (paralelo ao eixo de rotação da Terra).
    \item \textbf{Direção Principal ($+\hat{I}$):} Aponta para o Ponto Vernal ($\gamma$), ou Equinócio de Primavera. Esta direção é definida pela interseção entre o plano equatorial da Terra e o plano da Eclíptica (plano da órbita da Terra ao redor do Sol).
    \item \textbf{Eixo ($+\hat{J}$):} Completa o sistema dextrogiro ($\hat{J} = \hat{K} \times \hat{I}$), estando a $90^\circ$ a leste do ponto vernal no plano equatorial.
\end{itemize}

\begin{figure}[ht]
	\caption{Definição do Sistema de Coordenadas Geocêntrico Inercial (ECI).}
	\vspace{3mm}
	\begin{center}
    	\includegraphics[width=\mylenfig]{docs/figuras/eci_definition.png}  
	\end{center}
	\vspace{4mm}
	\legenda{O sistema ECI é fixo em relação às estrelas distantes. O eixo $X$ aponta para o Ponto Vernal ($\gamma$), definido pela interseção entre o plano equatorial e a eclíptica. O eixo $Z$ coincide com o eixo de rotação da Terra, apontando para o Polo Norte Celeste, enquanto o eixo $Y$ completa a tríade cartesiana dextrogira no plano equatorial ($90^\circ$ a leste de $\gamma$).}
	\label{fig:eci_frame}
	\FONTE{Adaptada de \citeonline{Vallado2013}.}
\end{figure}

A orientação do eixo de rotação terrestre e, consequentemente, do plano equatorial, não é estática no espaço. Devido aos torques gravitacionais exercidos principalmente pelo Sol e pela Lua sobre o bojo equatorial da Terra, o planeta experimenta movimentos de precessão e nutação.

Esses fenômenos fazem com que a posição do Ponto Vernal ($\gamma$) e do Polo Celeste varie continuamente, o que introduziria forças fictícias caso o sistema de coordenadas acompanhasse a data atual (\textit{True of Date}). Para garantir a inercialidade necessária à integração das Leis de Newton, adota-se uma época de referência fixa para a orientação dos eixos. O padrão internacional vigente é o sistema J2000.0 (também denotado como EME2000 - \textit{Earth Mean Equator 2000}). Este referencial define os eixos $X, Y, Z$ com base no equador médio e equinócio médio da época Juliana 2451545.0 (correspondente a 1º de janeiro de 2000, às 12:00:00 TT - Tempo Terrestre).

É neste referencial que os vetores posição $\vec{r}$ e velocidade $\vec{v}$ são definidos para a integração das equações de movimento.

\subsection{O Problema de Dois Corpos}
\label{subsec:dois_corpos}

O modelo fundamental para a descrição da dinâmica orbital é o Problema de Dois Corpos. Esta formulação idealizada assume que a Terra e o satélite comportam-se como massas pontuais e que a única força atuante no sistema é a atração gravitacional mútua, desconsiderando quaisquer perturbações externas. A equação vetorial do movimento em um sistema de referência inercial (ECI) é expressa por \cite{Bate1971, Vallado2013}:

\begin{equation}
    \ddot{\vec{r}} = -\frac{\mu}{r^3}\vec{r}
    \label{eq:dois_corpos}
\end{equation}

onde $\vec{r}$ representa o vetor posição do satélite, $r$ a sua magnitude, e $\mu$ é o parâmetro gravitacional padrão da Terra ($\mu \approx 3,986 \times 10^{14} \, \text{m}^3/\text{s}^2$).

A natureza conservativa do campo gravitacional central implica a existência de integrais de movimento fundamentais. A primeira é o vetor momento angular específico, $\vec{h}$, definido como o produto vetorial entre a posição e a velocidade:

\begin{equation}
    \vec{h} = \vec{r} \times \vec{v} = \text{constante}
    \label{eq:momento_angular}
\end{equation}

A constância de $\vec{h}$ (tanto em magnitude quanto em direção) demonstra que o movimento orbital ocorre em um plano fixo no espaço inercial, perpendicular a este vetor. A segunda integral é a energia mecânica específica, $\mathcal{E}$, que permanece constante ao longo da trajetória:

\begin{equation}
    \mathcal{E} = \frac{v^2}{2} - \frac{\mu}{r} = \text{constante}
    \label{eq:energia_especifica}
\end{equation}

A integração da equação do movimento (\ref{eq:dois_corpos}) pode ser realizada através da introdução do vetor excentricidade, $\vec{e}$ (também conhecido como vetor de Laplace-Runge-Lenz). Este vetor adimensional é uma constante de movimento fundamental que surge da integração da equação diferencial e possui um significado geométrico preciso:

\begin{equation}
    \vec{e} = \frac{\vec{v} \times \vec{h}}{\mu} - \frac{\vec{r}}{r}
    \label{eq:vetor_excentricidade}
\end{equation}

Fisicamente, o vetor excentricidade define a orientação da órbita dentro do seu plano. Ele aponta invariavelmente do centro de atração (foco principal) em direção ao perigeu (ponto de maior aproximação). Sua magnitude, $e = |\vec{e}|$, corresponde à excentricidade da cônica, determinando a forma da trajetória (circular se $e=0$, elíptica se $0 < e < 1$, etc.).

O produto escalar entre o vetor posição $\vec{r}$ e este vetor constante $\vec{e}$ permite deduzir a geometria da órbita sem a necessidade de integração temporal direta. Sabendo que o ângulo entre $\vec{r}$ e $\vec{e}$ é a anomalia verdadeira $\nu$, tem-se $\vec{r} \cdot \vec{e} = r e \cos(\nu)$. Aplicando as identidades vetoriais apropriadas à equação (\ref{eq:vetor_excentricidade}), chega-se à equação geral das cônicas em coordenadas polares:

\begin{equation}
    r(\nu) = \frac{p}{1 + e \cos(\nu)}
    \label{eq:conica_geral}
\end{equation}

O parâmetro $p$ (\textit{semi-latus rectum}) determina a largura da cônica medida no foco, sendo derivado diretamente do momento angular ($p = h^2/\mu$).

Para definir a escala global da órbita, introduz-se o semieixo maior ($a$). Geometricamente, em uma órbita elíptica, $a$ representa a média aritmética entre as distâncias de perigeu e apogeu. Fisicamente, este parâmetro está intrinsecamente ligado à energia mecânica específica ($\mathcal{E}$) do satélite: quanto maior o semieixo maior, maior é a energia da órbita.

As relações fundamentais que conectam as integrais de movimento ($\mathcal{E}, h$) aos parâmetros geométricos da forma ($e$) e tamanho ($a$) são dadas por:

\begin{equation}
    a = -\frac{\mu}{2\mathcal{E}}, \quad e = \sqrt{1 + \frac{2\mathcal{E}h^2}{\mu^2}}
\end{equation}

Desta forma, para o caso de satélites em órbitas fechadas (onde $\mathcal{E} < 0$), o semieixo maior assume valores positivos e finitos, definindo o tamanho da elipse, enquanto a excentricidade ($0 \le e < 1$) define o seu achatamento.

Para o caso específico de satélites em órbitas fechadas ($\mathcal{E} < 0$, resultando em $0 \le e < 1$), a trajetória é uma elipse. O estado do satélite no espaço tridimensional é univocamente definido por seis constantes de integração, os Elementos Orbitais Clássicos (COEs), ilustrados na Figura \ref{fig:elementos_orbitais}:

\begin{itemize}
    \item \textbf{Tamanho e Forma:} O semieixo maior ($a$) define a dimensão da órbita e sua energia específica, enquanto a excentricidade ($e$) descreve o grau de achatamento da elipse.
    \item \textbf{Orientação no Espaço:} A inclinação ($i$) define o ângulo entre o plano orbital (vetor $\vec{h}$) e o plano equatorial da Terra. A ascensão reta do nodo ascendente ($\Omega$) define a orientação do plano orbital em relação ao Ponto Vernal (eixo $X$ inercial). O argumento do perigeu ($\omega$) define a orientação da elipse dentro do próprio plano orbital.
    \item \textbf{Posição do Satélite:} A anomalia verdadeira ($\nu$) localiza o satélite em sua trajetória num instante específico, medido angularmente a partir do perigeu.
\end{itemize}

\begin{figure}[!h]
    \caption{Representação geométrica dos Elementos Orbitais Clássicos (COEs) no sistema de referência inercial.}
    \vspace{6mm}
    \begin{center}
        \includegraphics[width=\mylenfig]{docs/figuras/elementos_orbitais_classicos.png}
    \end{center}
    \vspace{4mm}
    \legenda{O diagrama ilustra os seis parâmetros que definem o estado do satélite. O tamanho da órbita é determinado pelo semieixo maior ($a$). A orientação do plano orbital é fixada pela inclinação ($i$) e pela ascensão reta do nodo ascendente ($\Omega$). A orientação da elipse no plano é dada pelo argumento do perigeu ($\omega$), definido como o ângulo entre a linha dos nodos e o vetor excentricidade ($\vec{e}$), que aponta para o perigeu. A posição instantânea é dada pela anomalia verdadeira ($\nu$).}
    \label{fig:elementos_orbitais}
    \FONTE{Produção do autor.}
\end{figure}

\subsection{Conversão de Elementos para Vetores de Estado}
\label{subsec:coe_to_rv}

Enquanto os Elementos Orbitais Clássicos oferecem uma visão geométrica da órbita, as simulações numéricas operam no espaço de estados cartesiano. A conversão requer o relacionamento entre a geometria da elipse e a evolução temporal do satélite. Como a velocidade angular do satélite varia ao longo da órbita (Segunda Lei de Kepler), a anomalia verdadeira ($\nu$) não é uma função linear do tempo. Para contornar essa dificuldade, introduzem-se duas grandezas auxiliares:

\begin{enumerate}
    \item \textbf{Anomalia Excêntrica ($E$):} Um ângulo auxiliar medido no centro da elipse, obtido pela projeção da posição do satélite sobre uma circunferência circunscrita de raio $a$.
    \item \textbf{Anomalia Média ($M$):} Um ângulo fictício que cresce linearmente com o tempo, representando a posição angular que o satélite teria se orbitasse em um círculo com velocidade constante equivalente ao período real da órbita.
\end{enumerate}

O algoritmo de conversão segue uma sequência lógica: propagação temporal, resolução geométrica e rotação espacial.

A dinâmica temporal inicia-se pelo cálculo do movimento médio ($n$), que define a velocidade angular média do satélite:

\begin{equation}
    n = \sqrt{\frac{\mu}{a^3}}
    \label{eq:movimento_medio}
\end{equation}

Considerando um instante inicial $t_0$ onde a anomalia média é $M_0$, a anomalia média $M$ no instante de interesse $t$ é obtida linearmente:

\begin{equation}
    M = M_0 + n(t - t_0)
    \label{eq:propagacao_media}
\end{equation}

A relação fundamental entre a posição temporal ($M$) e a posição geométrica auxiliar ($E$) é dada pela \textbf{Equação de Kepler}:

\begin{equation}
    M = E - e \sin(E)
    \label{eq:kepler}
\end{equation}

Como esta é uma equação transcendental, $E$ deve ser obtida por métodos numéricos iterativos. Com $E$ determinado, calcula-se a anomalia verdadeira $\nu$ através das relações geométricas da elipse:

\begin{equation}
    \tan\left(\frac{\nu}{2}\right) = \sqrt{\frac{1+e}{1-e}} \tan\left(\frac{E}{2}\right)
    \label{eq:anomalia_verdadeira_E}
\end{equation}

O segundo passo consiste em determinar os vetores no sistema auxiliar perifocal ($PQW$). Este referencial é fixo no plano da órbita, com o eixo $P$ apontando para o perigeu. Primeiramente, define-se o parâmetro semi-latus rectum ($p$):

\begin{equation}
    p = a(1-e^2)
\end{equation}

Neste sistema, as componentes vetoriais são dadas por:

\begin{equation}
    \vec{r}_{PQW} = \begin{bmatrix} r \cos \nu \\ r \sin \nu \\ 0 \end{bmatrix}
    , \quad
    \vec{v}_{PQW} = \sqrt{\frac{\mu}{p}} \begin{bmatrix} -\sin \nu \\ e + \cos \nu \\ 0 \end{bmatrix}
    \label{eq:vetores_pqw}
\end{equation}

Finalmente, os vetores são transformados do sistema perifocal para o sistema inercial geocêntrico (ECI) através de uma sequência de rotações de Euler envolvendo os ângulos de orientação da órbita ($\Omega, i, \omega$). A matriz de rotação resultante $[R]$ é dada por:

\begin{equation}
    [R] = \begin{bmatrix}
    c_{\Omega} c_{\omega} - s_{\Omega} s_{\omega} c_i & -c_{\Omega} s_{\omega} - s_{\Omega} c_{\omega} c_i & s_{\Omega} s_i \\
    s_{\Omega} c_{\omega} + c_{\Omega} s_{\omega} c_i & -s_{\Omega} s_{\omega} + c_{\Omega} c_{\omega} c_i & -c_{\Omega} s_i \\
    s_{\omega} s_i & c_{\omega} s_i & c_i
    \end{bmatrix}
    \label{eq:matriz_rotacao}
\end{equation}

Onde $c_{(\cdot)}$ e $s_{(\cdot)}$ representam cosseno e seno, respectivamente. Assim, os vetores de estado finais no sistema inercial são obtidos pelas operações matriciais:

\begin{equation}
    \vec{r}_{ECI} = [R] \vec{r}_{PQW}, \quad \vec{v}_{ECI} = [R] \vec{v}_{PQW}
    \label{eq:conversao_final}
\end{equation}

Embora o Problema de Dois Corpos forneça a solução kepleriana fundamental, ele baseia-se em hipóteses simplificadoras — como a esfericidade perfeita da Terra e a ausência de forças não-gravitacionais — que não representam a realidade completa do ambiente orbital. Para aplicações de alta precisão, como a análise de risco de colisão, é necessário recorrer a modelos perturbados.

\subsection{Perturbações Orbitais}
\label{subsec:perturbacoes}

O modelo de Dois Corpos fornece a solução para uma órbita ideal, definida por constantes de movimento. No entanto, o ambiente espacial real impõe forças adicionais. Para tratar matematicamente este problema sem abandonar a intuição geométrica das cônicas, utiliza-se o conceito de Elementos Osculadores.

Neste formalismo, assume-se que, em qualquer instante $t$, o satélite está sobre uma órbita kepleriana tangencial à trajetória real. Porém, devido às perturbações, os parâmetros que definem essa elipse ($a, e, i, \Omega, \omega, M$) variam continuamente com o tempo.

\vspace{0.5cm}
\noindent \textbf{Equações Planetárias de Lagrange}

A evolução temporal dos elementos orbitais sob a influência de forças conservativas (como o campo gravitacional não-esférico) é descrita pelas Equações Planetárias de Lagrange. Estas equações relacionam a taxa de variação dos elementos ($\dot{\alpha}_i$) com as derivadas parciais de uma Função Perturbadora ($R$).

A função $R$ representa o potencial gravitacional excedente ao termo central monopolo ($-\mu/r$). A forma geral das equações para as variações temporais é dada por \cite{Vallado2013, Kaula1966}:

\begin{equation}
    \begin{aligned}
        \frac{da}{dt} &= \frac{2}{na} \frac{\partial R}{\partial M} \\
        \frac{de}{dt} &= \frac{1-e^2}{na^2e} \frac{\partial R}{\partial M} - \frac{\sqrt{1-e^2}}{na^2e} \frac{\partial R}{\partial \omega} \\
        \frac{di}{dt} &= \frac{\cos i}{na^2\sqrt{1-e^2}\sin i} \frac{\partial R}{\partial \omega} - \frac{1}{na^2\sqrt{1-e^2}\sin i} \frac{\partial R}{\partial \Omega} \\
        \frac{d\Omega}{dt} &= \frac{1}{na^2\sqrt{1-e^2}\sin i} \frac{\partial R}{\partial i} \\
        \frac{d\omega}{dt} &= -\frac{\cos i}{na^2\sqrt{1-e^2}\sin i} \frac{\partial R}{\partial i} + \frac{\sqrt{1-e^2}}{na^2e} \frac{\partial R}{\partial e} \\
        \frac{dM}{dt} &= n - \frac{1-e^2}{na^2e} \frac{\partial R}{\partial e} - \frac{2}{na} \frac{\partial R}{\partial a}
    \end{aligned}
    \label{eq:lagrange_planetary}
\end{equation}

A solução analítica destas equações revela que as variações nos elementos orbitais podem ser classificadas em três categorias distintas, ilustradas na Figura \ref{fig:raan_osculador_vs_medio}:

\begin{enumerate}
    \item \textbf{Variações Seculares:} Alterações monotônicas e progressivas que crescem linearmente com o tempo (ex: a precessão dos nodos devido ao $J_2$). São as mais importantes para a análise de missão a longo prazo.
    \item \textbf{Variações de Longo Período:} Oscilações cíclicas com período superior ao período orbital do satélite (frequentemente relacionadas ao movimento de rotação do perigeu, com ciclos de semanas ou meses).
    \item \textbf{Variações de Curto Período:} Oscilações rápidas que ocorrem dentro de uma única revolução orbital (frequentemente dependentes da anomalia verdadeira ou média).
\end{enumerate}

Esta distinção é fundamental na propagação de órbitas: os elementos médios capturam apenas o comportamento secular, enquanto os elementos osculadores incluem todas as variações periódicas instantâneas.

\begin{figure}[ht]
    \caption{Comparação entre a evolução temporal da Ascensão Reta do Nodo Ascendente ($\Omega$) osculadora e média.}
    \vspace{3mm}
    \begin{center}
        \includegraphics[width=\mylenfig]{docs/figuras/mean_vs_osculation.png}  
    \end{center}
    \vspace{4mm}
    \legenda{A linha sólida azul exibe o valor osculador de $\Omega$, contendo a soma de todas as perturbações. A linha tracejada vermelha representa o valor médio, isolando a tendência secular conhecida como Precessão Nodal (ou Regressão dos Nodos), causada pelo achatamento terrestre ($J_2$). O gráfico abrange 100 períodos orbitais para evidenciar a sobreposição das oscilações de curto período (ruído de alta frequência) e longo período (ondas cíclicas) sobre a deriva secular constante.}
    \label{fig:raan_osculador_vs_medio}
    \FONTE{Produção do autor.}
\end{figure}

\vspace{0.5cm}
\noindent \textbf{Achatamento Terrestre ($J_2$)}

Aplicando o formalismo de Lagrange ao potencial gravitacional da Terra, é necessário considerar que o planeta não é uma esfera perfeita com distribuição de massa homogênea. Devido à sua rotação, a Terra apresenta um achatamento nos polos e um acúmulo de massa na região equatorial (bojo). Essa geometria faz com que o campo gravitacional desvie do modelo de massa pontual ideal ($U = \mu/r$).

Matematicamente, esse potencial não-esférico é modelado através de uma expansão em harmônicos esféricos. O termo dominante nessa expansão é o segundo harmônico zonal, denotado por $J_2$, que quantifica o achatamento oblato da Terra. O potencial perturbador associado a $J_2$ é dado por \cite{Vallado2013}:

\begin{equation}
    R_{J_2} = -\frac{\mu}{r} J_2 \left( \frac{R_\oplus}{r} \right)^2 \left( \frac{3}{2} \sin^2 \phi - \frac{1}{2} \right)
\end{equation}

\noindent onde $\mu$ é o parâmetro gravitacional da Terra, $R_\oplus$ é o raio equatorial terrestre e $\phi$ é a latitude geocêntrica do satélite. O coeficiente $J_2$ possui valor aproximado de $1,082 \times 10^{-3}$, sendo cerca de três ordens de grandeza superior aos demais harmônicos ($J_3, J_4, \dots$), o que justifica seu tratamento isolado em análises de primeira ordem.

Ao substituir $R_{J_2}$ nas Equações Planetárias de Lagrange e calcular a média sobre um período orbital, os termos osculatórios se anulam, restando as variações seculares (as tendências constantes vistas na Figura \ref{fig:raan_osculador_vs_medio}).

As principais perturbações seculares causadas pelo $J_2$ afetam a Ascensão Reta do Nodo Ascendente ($\Omega$), o Argumento do Perigeu ($\omega$) e a Anomalia Média ($M$). As taxas de variação média são dadas por:

\begin{equation}
    \dot{\bar{\Omega}}_{J_2} = - \frac{3}{2} n J_2 \left( \frac{R_\oplus}{p} \right)^2 \cos i
    \label{eq:j2_raan}
\end{equation}

\begin{equation}
    \dot{\bar{\omega}}_{J_2} = \frac{3}{4} n J_2 \left( \frac{R_\oplus}{p} \right)^2 \left( 4 - 5 \sin^2 i \right)
    \label{eq:j2_argp}
\end{equation}

\begin{equation}
    \dot{\bar{M}}_{J_2} = n \left[ 1 + \frac{3}{4} J_2 \left( \frac{R_\oplus}{p} \right)^2 \sqrt{1-e^2} \left( 2 - 3 \sin^2 i \right) \right]
    \label{eq:j2_mean_anomaly}
\end{equation}

\noindent onde:
\begin{itemize}
    \item $n = \sqrt{\mu/a^3}$ é o movimento médio kepleriano não-perturbado;
    \item $p = a(1 - e^2)$ é o \textit{semilatus rectum};
    \item $i$ é a inclinação orbital;
    \item $\sqrt{1-e^2}$ é um termo de acoplamento da excentricidade.
\end{itemize}

A Equação \ref{eq:j2_raan} descreve a Precessão dos Nodos. A direção dessa precessão depende do sinal de $\cos i$:

\begin{itemize}
    \item Para inclinações prógradas ($0^\circ \leq i < 90^\circ$), a precessão é negativa (o plano gira para Oeste).
    \item Para inclinações retrógradas ($90^\circ < i \leq 180^\circ$), a precessão torna-se positiva (o plano gira para Leste).
\end{itemize}

Esta dependência angular da precessão nodal constitui o fundamento físico das Órbitas Heliossíncronas (SSO). O objetivo destas missões é sincronizar a rotação do plano orbital com a taxa média de translação da Terra ao redor do Sol ($\dot{\bar{\Omega}} \approx 0,9856^\circ/\text{dia}$).

Como a taxa alvo é positiva, a análise da Equação \ref{eq:j2_raan} revela a necessidade de um termo $\cos i$ negativo. Consequentemente, tais órbitas requerem o regime retrógrado ($i > 90^\circ$). Na prática, o equilíbrio entre altitude e inclinação resulta em valores típicos entre $97^\circ$ e $99^\circ$ para missões de sensoriamento remoto, como as séries CBERS e Landsat.

Já a Equação \ref{eq:j2_argp} descreve a Rotação da Linha dos Apsides. Nota-se que para uma inclinação crítica de $i \approx 63,4^\circ$ (onde $4 - 5 \sin^2 i = 0$), a rotação do perigeu é nula. Esta propriedade é utilizada nas órbitas do tipo Molniya para manter o apogeu fixo sobre o hemisfério norte.

\vspace{0.5cm}
\noindent \textbf{Arrasto Atmosférico}

Diferentemente das perturbações gravitacionais, que são conservativas, o arrasto atmosférico atua como uma força dissipativa, removendo energia mecânica da órbita. Para satélites em Órbita Terrestre Baixa (LEO, altitude $< 1000$ km), esta é a perturbação dominante, sendo a principal responsável pelo decaimento orbital e pela determinação do tempo de vida útil da missão.

A força de arrasto surge da interação aerodinâmica entre a superfície do satélite e as partículas da alta atmosfera (termosfera e exosfera). A aceleração perturbadora resultante, $\mathbf{a}_{drag}$, é modelada classicamente por \cite{Montenbruck2000}:

\begin{equation}
    \mathbf{a}_{drag} = - \frac{1}{2} \rho \frac{C_D A}{m} v_{rel}^2 \frac{\mathbf{v}_{rel}}{|\mathbf{v}_{rel}|}
    \label{eq:drag_force}
\end{equation}

\noindent onde:
\begin{itemize}
    \item $\rho$ é a densidade atmosférica local;
    \item $C_D$ é o coeficiente de arrasto (adimensional), que depende da geometria do satélite e do regime de interação molecular (tipicamente $2,0 \leq C_D \leq 2,2$ para formas complexas);
    \item $A$ é a área de seção transversal projetada na direção do fluxo;
    \item $m$ é a massa do satélite;
    \item $\mathbf{v}_{rel}$ é o vetor velocidade do satélite relativo à atmosfera rotativa.
\end{itemize}

O termo $\frac{m}{C_D A}$ é frequentemente denominado Coeficiente Balístico ($B$). Um alto coeficiente balístico indica que o objeto tem grande inércia em relação à sua área de arrasto, sendo menos suscetível à desaceleração atmosférica.

A principal fonte de incerteza na Equação \ref{eq:drag_force} reside na modelagem da densidade atmosférica ($\rho$). A alta atmosfera é um meio dinâmico, onde a densidade não segue apenas um perfil de decaimento exponencial com a altitude, mas varia temporalmente em função da atividade solar e geomagnética. O aquecimento da termosfera, impulsionado pela radiação ultravioleta (monitorada pelo fluxo F10.7) e tempestades magnéticas, provoca a expansão das camadas atmosféricas. Esse fenômeno pode elevar a densidade local em uma altitude fixa em ordens de grandeza, alterando significativamente a força de arrasto prevista \cite{Jacchia1977}.

Do ponto de vista da mecânica orbital, o arrasto tem dois efeitos seculares principais descritos por \citeonline{KingHele1987}:
\begin{enumerate}
    \item \textbf{Decaimento do Semieixo Maior ($a$):} A perda de energia cinética reduz a altitude média da órbita, levando à reentrada eventual na atmosfera densa.
    \item \textbf{Circularização da Órbita:} O arrasto é mais intenso no perigeu (onde a densidade e a velocidade são máximas). O efeito de frenagem neste ponto reduz a altura do apogeu subsequente. Consequentemente, a excentricidade ($e$) tende a zero ao longo do tempo.
\end{enumerate}

Matematicamente, a evolução temporal dos elementos orbitais sob a ação de forças não-conservativas é descrita pelas Equações Gaussianas de Variação de Parâmetros (\textit{Gaussian VOP}). Diferentemente das equações planetárias de Lagrange, esta formulação permite tratar a perturbação diretamente através das componentes da força específica perturbadora decomposta no sistema local RSW ($F_r, F_s, F_w$):

\begin{itemize}
    \item $F_r$: Componente \textbf{Radial}, na direção do vetor posição;
    \item $F_s$: Componente \textbf{Transversal} (ou \textit{Along-track}), perpendicular ao raio vetor, no plano da órbita e no sentido do movimento;
    \item $F_w$: Componente \textbf{Normal} (ou \textit{Cross-track}), perpendicular ao plano orbital.
\end{itemize}

As taxas de variação dos elementos keplerianos clássicos são dadas pelo sistema \cite{Vallado2013}:

\begin{equation} \label{eq:gauss_a}
    \frac{da}{dt} = \frac{2a^2}{h} \left( e \sin \nu \, F_r + \frac{p}{r} \, F_s \right)
\end{equation}

\begin{equation} \label{eq:gauss_e}
    \frac{de}{dt} = \frac{1}{h} \left[ p \sin \nu \, F_r + \left( (p+r) \cos \nu + r e \right) F_s \right]
\end{equation}

\begin{equation} \label{eq:gauss_i}
    \frac{di}{dt} = \frac{r \cos u}{h} \, F_w
\end{equation}

\begin{equation} \label{eq:gauss_raan}
    \frac{d\Omega}{dt} = \frac{r \sin u}{h \sin i} \, F_w
\end{equation}

\begin{equation} \label{eq:gauss_omega}
    \frac{d\omega}{dt} = \frac{1}{h e} \left[ -p \cos \nu \, F_r + (p+r) \sin \nu \, F_s \right] - \frac{r \sin u \cos i}{h \sin i} \, F_w
\end{equation}

\begin{equation} \label{eq:gauss_M}
    \frac{dM}{dt} = n - \frac{b}{a h e} \left[ (p \cos \nu - 2re) F_r - (p+r) \sin \nu \, F_s \right]
\end{equation}

\noindent onde:
\begin{itemize}
    \item $u = \omega + \nu$ é o argumento da latitude;
    \item $h = \sqrt{\mu p}$ é o momento angular específico (magnitude do vetor $\mathbf{r} \times \mathbf{v}$);
    \item $p = a(1-e^2)$ é o \textit{semilatus rectum};
    \item $b = a\sqrt{1-e^2}$ é o semieixo menor da órbita;
    \item $n$ é o movimento médio do satélite ($\sqrt{\mu/a^3}$).
\end{itemize}

Aplicando este formalismo ao arrasto atmosférico, nota-se que a força atua predominantemente na direção oposta à velocidade. Para órbitas circulares ou pouco excêntricas, a componente Transversal ($F_s$) é dominante e negativa. Observando a Equação \ref{eq:gauss_a}, verifica-se que um valor negativo de $F_s$ resulta diretamente em $da/dt < 0$, confirmando matematicamente o decaimento orbital. Por outro lado, como o arrasto não possui componente significativa fora do plano orbital ($F_w \approx 0$), as equações \ref{eq:gauss_i} e \ref{eq:gauss_raan} indicam que a inclinação e a ascensão reta do nodo sofrem alterações desprezíveis devido à atmosfera.

Entretanto, embora as equações VOP sejam exatas, a incerteza estocástica na densidade atmosférica ($\rho$) torna inviável a obtenção de soluções analíticas precisas para longos períodos. Para aplicações de alta fidelidade, torna-se indispensável o uso de Métodos de Integração Numérica. As duas abordagens clássicas são o Método de Cowell, que integra diretamente a equação do movimento completa ($\ddot{\mathbf{r}} = -\frac{\mu}{r^3}\mathbf{r} + \mathbf{a}_{pert}$), e o Método de Encke, que integra apenas os desvios em relação a uma órbita de referência osculadora. Atualmente, devido ao aumento da capacidade computacional, o método de Cowell consolidou-se como o padrão para a propagação de órbitas LEO.

\vspace{0.5cm}
\noindent \textbf{Outras Perturbações}

Embora menos predominantes em LEO do que o $J_2$ e o arrasto, outras forças devem ser consideradas para propagação de alta fidelidade, como a atração gravitacional de terceiros corpos (Sol e Lua) e a Pressão de Radiação Solar (SRP), conforme discutido em \citeonline{Bate1971} e \citeonline{Vallado2013}.

\subsection{Propagadores Orbitais}
\label{subsec:propagadores}

A propagação orbital consiste no processo de determinar o estado futuro de um satélite ($\vec{r}, \vec{v}$) a partir de um estado inicial conhecido e de um modelo dinâmico. A escolha do algoritmo de propagação é ditada pelo compromisso entre a precisão requerida e o tempo computacional disponível. Na prática operacional e na análise de risco de colisão, distinguem-se duas categorias principais: os modelos analíticos e os modelos numéricos.


\noindent \textbf{Modelos Analíticos e o Padrão SGP4}

Os modelos analíticos buscam uma solução fechada para as equações do movimento, integrando as variações dos elementos orbitais em vez de integrar as acelerações cartesianas passo a passo. O padrão global para a catalogação e rastreio de objetos espaciais é o algoritmo SGP4 (\textit{Simplified General Perturbations 4}).

Desenvolvido pelo NORAD na década de 1970 e refinado posteriormente \cite{Hoots1980, Vallado2006_TLE}, o SGP4 foi projetado especificamente para processar os dados distribuídos no formato TLE (\textit{Two-Line Element}). É fundamental compreender que o TLE e o SGP4 formam um par indissociável: os elementos contidos em um TLE não são elementos keplerianos osculadores, mas sim elementos médios específicos que recuperam a posição correta do satélite quando processados pelas equações do SGP4, que reconstrói as oscilações de curto e longo período.

O modelo matemático do SGP4 considera o campo gravitacional da Terra através dos harmônicos zonais principais ($J_2, J_3, J_4$) e o arrasto atmosférico modelado por uma função de densidade simplificada. No contexto do arrasto, o TLE fornece o parâmetro $B^*$, que é um termo modificado do coeficiente balístico ($C_D A/m$) adaptado para o modelo de atmosfera do SGP4.

A principal vantagem desta abordagem é a velocidade computacional, permitindo a propagação de grandes catálogos em frações de segundo. No entanto, sua precisão é limitada, apresentando erros típicos da ordem de 1 km na época do TLE, degradando-se com o tempo de propagação (1 a 3 km de erro por dia) devido às simplificações na modelagem das perturbações. Para a triagem inicial de conjunções (\textit{screening}), o SGP4 é a ferramenta padrão; contudo, para o cálculo refinado da probabilidade de colisão e planejamento de manobras, sua precisão é insuficiente.


\noindent \textbf{Propagação Numérica de Alta Fidelidade}

Para aplicações críticas como a Otimização de Manobras de Desvio (CAMO), torna-se indispensável o uso de Métodos de Integração Numérica (conhecidos como \textit{Special Perturbations}). Diferentemente da abordagem analítica, estes métodos integram diretamente a equação diferencial do movimento completa:

\begin{equation}
    \ddot{\vec{r}} = -\frac{\mu}{r^3}\vec{r} + \vec{a}_{pert}
\end{equation}

Onde $\vec{a}_{pert}$ é o somatório vetorial de todas as acelerações perturbadoras modeladas com alta fidelidade, incluindo o geo-potencial de alta ordem (harmônicos esféricos e tesserais), o arrasto com atmosfera dinâmica, a atração gravitacional de terceiros corpos (Sol e Lua) e a pressão de radiação solar.

Existem duas abordagens clássicas para esta integração. O Método de Cowell integra diretamente os vetores de posição e velocidade totais. É o método mais direto de implementar e, com o avanço da aritmética de dupla precisão e da capacidade de processamento moderna, consolidou-se como o padrão industrial. Alternativamente, o Método de Encke integra apenas os desvios em relação a uma órbita de referência osculadora, o que era historicamente vantajoso para reduzir erros de arredondamento em computadores antigos, mas hoje é menos prevalente em simulações LEO genéricas.

A propagação numérica oferece precisão superior, mas com custo computacional elevado. Já o SGP4 é extremamente rápido, porém trabalha com elementos médios e menor precisão.

No contexto desta pesquisa, essa dicotomia é central. Uma das abordagens propostas investiga o uso híbrido: utilizar a velocidade do SGP4 acoplada às Equações Variacionais de Gauss para a evolução rápida da população no Algoritmo Genético, comparando posteriormente a qualidade dessas soluções com aquelas obtidas via propagação numérica de alta fidelidade. Essa análise comparativa visa determinar se o ganho de tempo computacional dos modelos analíticos compensa a perda de acurácia na estimativa do risco de colisão.

\section{Dinâmica e Manobras Orbitais}
\label{sec:dinamica_manobras}

A mitigação do risco de colisão exige a alteração ativa da trajetória do satélite. Nesta seção, definem-se os modelos matemáticos utilizados para a aplicação de impulsos de controle e para o planejamento de manobras de retorno à órbita nominal.

\subsection{Manobras Impulsivas}
\label{subsec:manobras_impulsivas}

No contexto da otimização de trajetórias, a modelagem mais comum assume a hipótese de queima impulsiva (\textit{Impulsive Burn}). Conforme descrito por \citeonline{Curtis2013}, nesta aproximação considera-se que o tempo de funcionamento dos propulsores ($\Delta t_{burn}$) é desprezável em comparação ao período orbital, de modo que a mudança de velocidade ocorre instantaneamente sem alteração na posição do veículo.

O vetor velocidade pós-manobra ($\vec{v}_{f}$) é dado por:
\begin{equation}
    \vec{v}_{f} = \vec{v}_{i} + \Delta \vec{v}
\end{equation}
\noindent onde $\vec{v}_{i}$ é a velocidade inercial imediatamente antes da queima e $\Delta \vec{v}$ é o vetor impulso fornecido pelo sistema de propulsão.

A magnitude $|\Delta \vec{v}|$ está diretamente relacionada ao consumo de propelente (``custo'' da manobra). A relação fundamental é dada pela Equação de Tsiolkovsky \cite{Sutton2016, Vallado2013}:
\begin{equation}
    \Delta v = I_{sp} g_0 \ln \left( \frac{m_{inicial}}{m_{final}} \right)
\end{equation}
\noindent onde $I_{sp}$ é o impulso específico do motor e $g_0$ a gravidade padrão.

Geralmente, o vetor $\Delta \vec{v}$ é decomposto no sistema local RSW (Radial, Transversal, Normal). Para manobras de evasão de colisão em LEO, a componente Transversal (tangencial à órbita) é a mais eficiente energeticamente para alterar o semieixo maior e, consequentemente, modificar o período orbital \cite{Vallado2013}. Esta alteração resulta em uma mudança no tempo de chegada ao ponto de encontro (estratégia de \textit{phasing}), aumentando a separação entre os objetos no instante de maior aproximação.

\subsection{Dinâmica do Movimento Relativo}
\label{subsec:movimento_relativo}

Em cenários operacionais complexos, não basta apenas evitar a colisão; é frequentemente necessário garantir que o satélite retorne à sua posição orbital nominal (ou \textit{slot} operacional) após o evento de risco. Para calcular o impulso necessário para este retorno, utiliza-se a dinâmica do movimento relativo.

O referencial padrão para descrever este movimento é o sistema LVLH (\textit{Local Vertical, Local Horizontal}). Neste trabalho, o referencial LVLH é definido em alinhamento direto com o sistema RSW anteriormente apresentado, onde o vetor de estado relativo $\delta \mathbf{x} = [x, y, z, \dot{x}, \dot{y}, \dot{z}]^T$ segue a convenção:

\begin{itemize}
    \item \textbf{Eixo $x$ (Radial):} Direção radial (R), apontando do centro da Terra para o satélite (Zenith).
    \item \textbf{Eixo $y$ (Transversal):} Direção do movimento (S), tangente à órbita no plano orbital.
    \item \textbf{Eixo $z$ (Normal):} Direção normal ao plano orbital (W), completando a tríade dextrogira.
\end{itemize}

Com esta definição, assume-se a hipótese simplificadora de que a excentricidade da órbita de referência é nula ou desprezável ($e \approx 0$). Esta aproximação é amplamente aceita na literatura para aplicações em LEO, onde a maioria das órbitas são quase-circulares. Sob estas condições, a dinâmica relativa linearizada é descrita pelas equações de Hill-Clohessy-Wiltshire (HCW) \cite{Clohessy1960, Vallado2013}:

\begin{equation} \label{eq:hcw_x}
    \ddot{x} - 2n\dot{y} - 3n^2x = f_x
\end{equation}
\begin{equation} \label{eq:hcw_y}
    \ddot{y} + 2n\dot{x} = f_y
\end{equation}
\begin{equation} \label{eq:hcw_z}
    \ddot{z} + n^2z = f_z
\end{equation}

\noindent onde $n = \sqrt{\mu/a^3}$ é o movimento médio da órbita de referência e $f_i$ são as componentes da aceleração de controle. Vale ressaltar que, para órbitas com excentricidade elevada ($e > 0.1$), estas equações perdem precisão, exigindo o uso das equações de Tschauner-Hempel, o que foge ao escopo das missões LEO típicas abordadas neste trabalho.

A solução analítica destas equações lineares fornece a Matriz de Transição de Estado (STM, $\mathbf{\Phi}(t, t_0)$), que mapeia o estado relativo atual para um estado futuro. Esta propriedade permite a aplicação direta da teoria de manobras impulsivas no referencial relativo.

Para o cálculo da manobra de retorno, formula-se o problema como um sistema linear. Dado o estado relativo pós-evasão em $t_0$ ($\delta \mathbf{r}_0$) e o objetivo de alcançar a origem do sistema ($\delta \mathbf{r}_f = \mathbf{0}$) em um tempo fixo $\Delta t$, é possível inverter os sub-blocos da matriz $\mathbf{\Phi}$ para determinar o vetor impulso de velocidade instantâneo ($\Delta \vec{v}$) necessário:

\begin{equation}
    \Delta \vec{v}_{req} = \mathbf{\Phi}_{rv}^{-1} (\Delta t) \left[ \delta \mathbf{r}_f - \mathbf{\Phi}_{rr}(\Delta t) \delta \mathbf{r}_0 \right] - \delta \vec{v}_0
\end{equation}

\noindent onde $\mathbf{\Phi}_{rr}$ e $\mathbf{\Phi}_{rv}$ são as submatrizes de partição da STM que relacionam, respectivamente, posição-posição e velocidade-posição. Esta abordagem, conhecida como \textit{Impulsive Rendezvous} ou Problema de Targeting, permite calcular analiticamente o custo de propelente para restaurar a órbita nominal após a manobra de evasão \cite{Alfriend2010, Vallado2013}.

\section{Análise de Conjunções e Risco de Colisão}
\label{sec:analise_conjuncao_risco}

A eficácia de uma manobra de evasão depende da capacidade de prever a trajetória futura dos objetos. Contudo, a determinação exata da posição de um satélite está sujeita a incertezas. Esta seção aborda a natureza estocástica da predição orbital e define as métricas quantitativas utilizadas para traduzir essa incerteza em critérios de decisão.

\subsection{Incertezas na Predição Orbital}
\label{subsec:incertezas_predicao}

Estimativas de estado orbital possuem erros inerentes. Conforme classificado por \citeonline{Klinkrad2006}, as incertezas originam-se de três fontes fundamentais: erros na determinação da órbita inicial (precisão dos sensores), erros nos modelos de força (densidade atmosférica e geopotencial) e erros numéricos de integração.

Em cenários de alta fidelidade, essa incerteza é representada matematicamente pela matriz de covariância de estado ($\mathbf{P}$), uma matriz simétrica $6 \times 6$ que correlaciona os desvios de posição e velocidade em relação à média:

\begin{equation}
    \mathbf{P} = \begin{bmatrix}
    \mathbf{C}_{rr} & \mathbf{C}_{rv} \\
    \mathbf{C}_{vr} & \mathbf{C}_{vv}
    \end{bmatrix}
\end{equation}

A evolução desta incerteza é regida pela dinâmica orbital. Pequenos erros iniciais de velocidade ($\mathbf{C}_{vv}$) propagam-se resultando em grandes erros de posição ($\mathbf{C}_{rr}$). Devido à correlação direta entre o semieixo maior e o período orbital, erros na estimativa de energia manifestam-se como um desvio cumulativo na direção do movimento do satélite (eixo $S$ do sistema RSW). Como consequência, o volume de incerteza tende a se alongar significativamente ao longo da trajetória orbital.

\subsection{Métricas de Risco de Colisão}
\label{subsec:metricas_risco}

Para fundamentar a decisão de manobra, a geometria do encontro e as incertezas devem ser expressas por uma métrica de risco objetiva. Duas métricas principais são utilizadas na astrodinâmica:

\textbf{\textit{Miss Distance}}: Métrica puramente geométrica, definida como a distância euclidiana mínima entre os centros de massa dos dois objetos no instante de maior aproximação (TCA). Embora intuitiva, possui limitações por ignorar a forma e a orientação da incerteza. Uma \textit{Miss Distance} de 1 km pode ser segura ou perigosa dependendo da magnitude da covariância de erro.

\textbf{Probabilidade de Colisão ($P_c$)}: Métrica estocástica que pondera a geometria do encontro pela incerteza associada. Representa a probabilidade de os objetos ocuparem o mesmo volume físico simultaneamente. Ao considerar a matriz de covariância combinada, a $P_c$ é a métrica de referência em operações espaciais, permitindo distinguir aproximações geometricamente próximas mas estatisticamente seguras de encontros com risco real de impacto.

\subsection{Estratégia de Seleção da Métrica de Risco}
\label{subsec:selecao_metricas}

A definição de qual das métricas de risco apresentadas acima será adotada como função objetivo neste trabalho não é arbitrária. A escolha deve ser estritamente coerente com a fidelidade dos dados de navegação disponíveis em cada cenário simulado:

\begin{itemize}
    \item \textbf{Cenários Analíticos (TLE/SGP4):} Como os dados TLE não possuem matriz de covariância formal, derivar probabilidades estatísticas resultaria em valores pouco confiáveis \cite{Klinkrad2006}. Portanto, para os cenários com SGP4, a métrica de risco selecionada é a \textbf{\textit{Miss Distance}}. A otimização busca maximizar a separação geométrica, priorizando a componente radial ($R$). Esta escolha justifica-se pois a incerteza do TLE na direção radial é historicamente a menor ($\sigma_R \ll \sigma_S$), permitindo garantir a segurança do encontro independentemente dos grandes erros de cronometragem (fase) associados à direção do movimento.
    
    \item \textbf{Cenários Numéricos:} Ao utilizar efemérides com propagação numérica da matriz de covariância $6 \times 6$, tem-se acesso à descrição estatística completa do erro. Nestes casos, a métrica de risco mandatória é a \textbf{Probabilidade de Colisão ($P_c$)}. O algoritmo de otimização atuará para minimizar o risco estocástico, explorando as regiões de baixa densidade de probabilidade, o que constitui uma abordagem superior à simples separação por distância.
\end{itemize}

\subsection{Identificação de Eventos de Conjunção}
\label{subsec:identificacao_conjuncoes}

A análise de risco de colisão impõe um desafio de complexidade computacional significativo. O catálogo público contém dezenas de milhares de objetos, tornando inviável a comparação numérica precisa da órbita do satélite operacional com todos os objetos catalogados a cada passo de integração. Para solucionar este problema, \citeonline{Klinkrad2006} descreve o método ``\textit{Smart Sieve}'', um algoritmo de filtragem hierárquica que elimina progressivamente pares de objetos não perigosos, reduzindo o espaço de busca antes do cálculo refinado do TCA e da métrica de risco (\textit{miss distance} ou probabilidade de colisão).

O processo inicia-se com o filtro radial (apogeu/perigeu), que é puramente geométrico e independe do tempo. Descarta-se pares cujas camadas orbitais não se interceptam fisicamente através da comparação dos raios de perigeu e apogeu:

\begin{equation}
    r_p = a(1-e)
\end{equation}
\begin{equation}
    r_a = a(1+e)
\end{equation}

Onde $a$ representa o semieixo maior e $e$ a excentricidade da órbita. O par é ignorado se a condição de exclusão abaixo for satisfeita, onde $\delta_r$ é uma margem de segurança para incertezas orbitais:

\begin{equation}
    (r_{p,1} - \delta_r) > (r_{a,2} + \delta_r) \quad \text{ou} \quad (r_{p,2} - \delta_r) > (r_{a,1} + \delta_r)
\end{equation}

Para os pares remanescentes, executa-se a triagem temporal discreta. Esta etapa utiliza filtros sequenciais com custos computacionais distintos para minimizar o processamento. As órbitas são propagadas em passos discretos $\Delta t$ e, para cada instante $t_k$, define-se o vetor posição relativa entre o objeto secundário e o primário:

\begin{equation}
    \Delta \vec{r}(t_k) = \vec{r}_2(t_k) - \vec{r}_1(t_k) = [x_{rel}, y_{rel}, z_{rel}]^T
\end{equation}

E o seu respectivo módulo, que representa a distância escalar entre os objetos:

\begin{equation}
    \rho(t_k) = \sqrt{x_{rel}^2 + y_{rel}^2 + z_{rel}^2}
\end{equation}

O primeiro filtro aplicado é o de menor custo computacional. Ele utiliza um raio de segurança ajustado ($R_{c,1}$) que depende da velocidade de escape da Terra ($v_e = \sqrt{2\mu/r}$). Para garantir que o filtro seja conservador, assume-se que a velocidade relativa máxima de encontro seja o dobro da velocidade de escape ($2v_e$). Contudo, como a análise de cada ponto $t_k$ deve cobrir apenas metade do intervalo de propagação para frente e para trás ($\Delta t / 2$), os fatores se cancelam, resultando na expressão:

\begin{equation}
    R_{c,1} = R_c + v_e \Delta t
\end{equation}

Nesta etapa, avalia-se inicialmente cada componente do vetor posição relativa de forma individual. O ponto $t_k$ é descartado imediatamente se qualquer uma das coordenadas satisfizer a condição de descarte rápido, o que evita o cálculo desnecessário da norma do vetor:

\begin{equation}
    |x_{rel}| > R_{c,1} \quad \text{ou} \quad |y_{rel}| > R_{c,1} \quad \text{ou} \quad |z_{rel}| > R_{c,1}
\end{equation}

Caso o ponto sobreviva a essa triagem inicial pelas componentes, realiza-se então a verificação do módulo do vetor distância relativa. O ponto é descartado se a distância euclidiana for superior ao raio de segurança ajustado:

\begin{equation}
    \rho(t_k) > R_{c,1}
\end{equation}

Se o ponto passar por essa sequência, aplica-se o filtro de curvatura $R_{c,2}$. Este estágio incorpora a aceleração da gravidade da Terra na superfície $g$ (análogo ao utilizar a velocidade de escape no cálculo de $R_{c,1}$) e a geometria curvilínea da órbita para estimar um raio de segurança mais preciso, isolando a componente da distância perpendicular ao vetor velocidade relativa $\Delta \vec{v}$:

\begin{equation}
    R_{c,2} = R_c + g(\Delta t)^2 < \sqrt{\rho^2 - \left(\frac{\Delta \vec{r} \cdot \Delta \vec{v}}{|\Delta \vec{v}|}\right)^2}
\end{equation}

A triagem é finalizada pela aplicação de um raio de segurança ainda mais preciso, $R_{c,3}$, que substitui o limite superior velocidade de escape pela velocidade relativa efetiva medida entre os objetos, otimizando o descarte antes da etapa de convergência numérica:

\begin{equation}
    R_{c,3} = R_{c,2} + \frac{1}{2} \left| \frac{\Delta \vec{v} \cdot \Delta \vec{r}}{\rho} \right| \Delta t
\end{equation}

Uma vez identificado um evento potencial ($\rho(t_k) < R_{c,3}$), um algoritmo determinístico de cálculo de mínimo local é utilizado para determinar o instante de máxima aproximação (TCA):

\begin{equation}
    \frac{\Delta \vec{v}_{TCA} \cdot \Delta \vec{r}_{TCA}}{|\Delta \vec{r}_{TCA}|} = \dot{\rho}_{TCA} = 0
\end{equation}

Em seguida, uma avaliação é realizada com base no vetor posição relativa no TCA no sistema RSW. Para converter os vetores de estado do sistema inercial (ECI) para o referencial local da órbita, realiza-se uma rotação de matrizes baseada nos vetores unitários que definem as direções Radial ($\hat{U}$), Along-track ($\hat{V}$) e Cross-track ($\hat{W}$):

\begin{equation}
    \hat{U} = \frac{\vec{r}}{|\vec{r}|}, \quad \hat{W} = \frac{\vec{r} \times \vec{v}}{|\vec{r} \times \vec{v}|}, \quad \hat{V} = \hat{W} \times \hat{U}
\end{equation}

A matriz de rotação $\mathbf{R}_{RSW}$ é construída a partir desses vetores, permitindo a transformação do vetor posição relativa $\Delta \vec{r}_{ECI}$ para o sistema de coordenadas intrínseco:

\begin{equation}
    \mathbf{R}_{RSW} = \begin{pmatrix} 
    \hat{U}_x & \hat{U}_y & \hat{U}_z \\ 
    \hat{V}_x & \hat{V}_y & \hat{V}_z \\ 
    \hat{W}_x & \hat{W}_y & \hat{W}_z 
    \end{pmatrix}
\end{equation}

O vetor posição relativa no referencial de colisão é então obtido por:

\begin{equation}
    \Delta \vec{r}_{RSW} = \mathbf{R}_{RSW} \Delta \vec{r}_{ECI} = [r_U, r_V, r_W]^T
\end{equation}



Esta transformação é fundamental para a validação final da conjunção, uma vez que os limites de segurança operacionais são definidos por um volume elipsoidal orientado nestes eixos. O evento é classificado como uma conjunção válida se o vetor posição relativa violar o volume de proteção, satisfazendo a norma quadrática baseada nos semi-eixos de segurança $R_{c,U}, R_{c,V}$ e $R_{c,W}$:

\begin{equation}
    k_c^2 = \left(\frac{r_U}{R_{c,U}}\right)^2 + \left(\frac{r_V}{R_{c,V}}\right)^2 + \left(\frac{r_W}{R_{c,W}}\right)^2 \leq 1
\end{equation}

Onde o parâmetro $k_c$ determina a penetração do objeto secundário no elipsoide de segurança do satélite operacional. Caso $k_c^2 \leq 1$, o evento é registrado como uma conjunção crítica e prossegue-se para o cálculo da probabilidade de colisão.

\subsection{Geometria do Encontro e o Plano B}
\label{subsec:plano_b}

Uma vez confirmada a conjunção crítica no referencial RSW, a análise de risco prossegue para a definição da geometria do encontro no instante de máxima aproximação ($t_{TCA}$). O cálculo da probabilidade de colisão é significativamente simplificado ao projetar as incertezas posicionais de ambos os objetos em um plano bidimensional perpendicular à velocidade relativa, conhecido como Plano de Colisão ou Plano B (\textit{B-Plane}).



Neste referencial, assume-se que as trajetórias dos objetos são lineares durante o curto intervalo do encontro, e que a velocidade relativa $\Delta \vec{v}$ permanece constante. O Plano B é definido de forma que o vetor velocidade relativa seja normal à sua superfície. Para construir este sistema de coordenadas local $(\hat{x}, \hat{y}, \hat{z})$, utiliza-se o vetor unitário na direção da velocidade relativa como o eixo secundário:

\begin{equation}
    \hat{z}_B = \frac{\Delta \vec{v}}{|\Delta \vec{v}|}
\end{equation}

Os eixos contidos no plano, $\hat{x}_B$ e $\hat{y}_B$, são escolhidos para representar a projeção das incertezas. Tipicamente, utiliza-se a projeção do vetor posição relativa ou a orientação das matrizes de covariância para definir esses eixos. A principal vantagem desta representação é reduzir a integração da função de densidade de probabilidade (PDF) 3D para uma integral dupla sobre uma área circular (ou elipsoidal) que representa a soma dos raios físicos dos objetos:

\begin{equation}
    P_c = \frac{1}{2\pi \sigma_x \sigma_y \sqrt{1-\rho_{xy}^2}} \iint_{A} \exp \left[ -\frac{1}{2} Q(x,y) \right] dx dy
\end{equation}

Onde $A$ é a área de seção transversal combinada e $Q(x,y)$ descreve a elipse de incerteza projetada no plano.

O Plano B é um referencial bidimensional centrado no objeto primário, construído de forma que o eixo normal ao plano coincida com a direção do vetor velocidade relativa $\Delta \vec{v}$ no instante de máxima aproximação. A principal premissa para a utilização deste plano é a de que, dada a magnitude das velocidades orbitais, o tempo de passagem de um objeto pelo volume de incerteza do outro é extremamente curto, permitindo tratar o movimento relativo como retilíneo e uniforme durante o encontro.



Para a construção da base ortonormal $\{\hat{x}_B, \hat{y}_B, \hat{z}_B\}$ do Plano B, define-se primeiramente o eixo unitário $\hat{z}_B$, paralelo à velocidade relativa:

\begin{equation}
    \hat{z}_B = \frac{\Delta \vec{v}}{|\Delta \vec{v}|}
\end{equation}

Os eixos contidos no plano de colisão, $\hat{x}_B$ e $\hat{y}_B$, devem ser ortogonais entre si e a $\hat{z}_B$. Embora existam diferentes convenções para a orientação destes eixos, uma escolha comum na literatura de dinâmica espacial consiste em projetar os componentes da posição relativa e das covariâncias de modo a minimizar o acoplamento entre as variáveis. Uma das formas de definir os eixos do plano é utilizar a direção do momento angular orbital ou a projeção do vetor posição inercial, garantindo a ortogonalidade via produto vetorial:

\begin{equation}
    \hat{x}_B = \frac{\vec{v}_1 \times \hat{z}_B}{|\vec{v}_1 \times \hat{z}_B|}
\end{equation}
\begin{equation}
    \hat{y}_B = \hat{z}_B \times \hat{x}_B
\end{equation}

Nesta configuração, o vetor distância de máxima aproximação $\Delta \vec{r}_{TCA}$ reside inteiramente no plano $x_B y_B$, uma vez que, por definição de $TCA$, o produto escalar $\Delta \vec{r} \cdot \Delta \vec{v} = 0$. Assim, a geometria 3D do encontro é reduzida a uma análise 2D, onde a distância de erro (\textit{miss distance}) $\rho$ é dada por:

\begin{equation}
    \rho = \sqrt{x_{B}^2 + y_{B}^2}
\end{equation}

O cálculo da probabilidade de colisão requer a combinação das incertezas posicionais de ambos os objetos. Sob a premissa de que os erros de estado do objeto primário ($\mathbf{C}_1$) e secundário ($\mathbf{C}_2$) são estatisticamente independentes, a matriz de covariância combinada $\mathbf{C}_c$ no referencial inercial é obtida pela soma das matrizes individuais:

\begin{equation}
    \mathbf{C}_c = \mathbf{C}_1 + \mathbf{C}_2
\end{equation}

Para projetar esta incerteza tridimensional diretamente no Plano B, define-se uma matriz de transformação $\mathbf{M}$ de dimensão $2 \times 3$. Esta matriz utiliza apenas os versores unitários $\hat{x}_B$ e $\hat{y}_B$, que definem o plano perpendicular à velocidade relativa:

\begin{equation}
    \mathbf{M} = \begin{pmatrix} 
    \hat{x}_{B,x} & \hat{x}_{B,y} & \hat{x}_{B,z} \\ 
    \hat{y}_{B,x} & \hat{y}_{B,y} & \hat{y}_{B,z} 
    \end{pmatrix} = \begin{bmatrix} \hat{x}_B^T \\ \hat{y}_B^T \end{bmatrix}
\end{equation}

Ao aplicar esta matriz à covariância combinada $\mathbf{C}_c$, obtém-se diretamente a matriz de covariância reduzida $\mathbf{C}_{2D}$ ($2 \times 2$), que representa a elipse de incerteza projetada no plano de encontro:

\begin{equation}
    \mathbf{C}_{2D} = \mathbf{M} \mathbf{C}_c \mathbf{M}^T = \begin{pmatrix} 
    \sigma_x^2 & \rho_{xy}\sigma_x\sigma_y \\ 
    \rho_{xy}\sigma_x\sigma_y & \sigma_y^2 
    \end{pmatrix}
\end{equation}

Esta abordagem simplifica o processo computacional, isolando apenas a dispersão estatística contida no plano onde o risco de colisão é avaliado. Como o movimento relativo é assumido como linear e a velocidade relativa é constante durante o curto intervalo do TCA, a componente da incerteza na direção do movimento ($\hat{z}_B$) é irrelevante para o cálculo da probabilidade de colisão instantânea. A densidade de probabilidade bidimensional resultante permite, então, que o problema seja resolvido através da integração sobre a área de seção transversal combinada dos objetos.

\subsection{Simplificação de Chan e Cálculo Analítico da Probabilidade}

O cálculo da probabilidade de colisão ($P_c$) fundamenta-se na integração da função densidade de probabilidade (PDF) gaussiana bidimensional sobre a área de seção transversal combinada dos objetos envolvidos. Conforme proposto por \citeonline{Chan2008}, a geometria complexa dos satélites é simplificada por um disco de raio efetivo $R_{obj}$ no plano de encontro, definido pela soma dos raios das esferas envolventes de ambos os corpos: $R_{obj} = R_1 + R_2$.

No Plano B, a PDF é centrada nas coordenadas de máxima aproximação $(x_m, y_m)$. Estes parâmetros representam a transposição do vetor posição relativa no instante de máxima aproximação ($\Delta \vec{r}_{TCA}$) para o referencial do plano. Dado que, no TCA, a condição de ortogonalidade $\Delta \vec{r}_{TCA} \cdot \Delta \vec{v}_{TCA} = 0$ é estritamente satisfeita, o vetor de separação mínima reside inteiramente no plano de colisão. Por conseguinte, sua representação bidimensional é obtida pela transformação linear baseada na matriz de projeção $\mathbf{M}$ ($2 \times 3$), constituída pelos versores unitários $\hat{x}_B$ e $\hat{y}_B$:

\begin{equation}
    \begin{Bmatrix} x_m \\ y_m \end{Bmatrix} = \mathbf{M} \Delta \vec{r}_{TCA}
\end{equation}

O módulo desse vetor projetado, $\rho = \sqrt{x_m^2 + y_m^2}$, define a \textit{miss distance}. A probabilidade de colisão é, então, expressa pela integral da PDF estruturada pela matriz de covariância projetada $\mathbf{C}_{2D}$:

\begin{equation}
    P_c = \frac{1}{2\pi \sqrt{|\mathbf{C}_{2D}|}} \iint_{A} \exp \left[ -\frac{1}{2} \begin{pmatrix} x \\ y \end{pmatrix}^T \mathbf{C}_{2D}^{-1} \begin{pmatrix} x \\ y \end{pmatrix} \right] dx dy
\end{equation}

Para viabilizar uma solução analítica e evitar a integração numérica, \citeonline{Chan2008} aplica uma operação de ``circularização''. O procedimento ocorre em duas etapas lógicas: primeiro, o sistema de coordenadas é rotacionado para alinhar-se com os eixos principais da elipse de incerteza (diagonalização da matriz $\mathbf{C}_{2D}$); em seguida, as dimensões físicas são reescalonadas, dividindo-se as coordenadas pelos seus respectivos desvios padrão ($\sigma_x, \sigma_y$).

Dessa forma, a distância de erro deixa de ser avaliada apenas em metros e passa a ser tratada como uma distância estatística adimensional. Esse processo garante que a probabilidade pondere corretamente a geometria do encontro: um erro de posição em uma direção de alta incerteza (grande $\sigma$) representa um risco estatístico menor do que a mesma distância em uma direção onde a posição é conhecida com alta precisão.

Sob a premissa de que o raio do objeto é pequeno comparado à magnitude das incertezas, a integral converge para uma expansão em série de potências de Rice \cite{Akella2000}:

\begin{equation}
    P_c = e^{-v/2} \sum_{m=0}^{\infty} \frac{(v/2)^m}{m!} \left[ 1 - e^{-u/2} \sum_{k=0}^{m} \frac{(u/2)^k}{k!} \right]
\end{equation}

Esta série é governada pelos parâmetros adimensionais $u$ e $v$, que sintetizam as características geométricas e estatísticas do encontro no espaço transformado:

\begin{equation}
    u = \frac{R_{obj}^2}{\sigma_x \sigma_y} \quad \text{e} \quad v = \frac{x_m^2}{\sigma_x^2} + \frac{y_m^2}{\sigma_y^2}
\end{equation}

Neste contexto, $v$ corresponde ao quadrado da distância de Mahalanobis, representando a separação estatística entre os objetos, enquanto $u$ define a relação entre a área física de risco e a dispersão da covariância. A formulação destaca-se pela rápida convergência, sendo suficiente para a maioria das aplicações operacionais de análise de risco.

\section{Algoritmos Genéticos}
\label{sec:algoritmos_geneticos}

A otimização de manobras orbitais é intrinsecamente caracterizada por um espaço de busca não-linear, multimodal e, frequentemente, não-convexo. Neste cenário, métodos clássicos baseados em gradiente, como o método de Newton-Raphson ou a Programação Quadrática Sequencial (SQP), apresentam limitações significativas, tendendo a convergir prematuramente para mínimos locais dependendo da estimativa inicial fornecida. Para superar tais barreiras e explorar o espaço de soluções de forma global, a literatura de astrodinâmica tem recorrido extensivamente a métodos heurísticos estocásticos, com destaque para os Algoritmos Genéticos (AG).

\subsection{Fundamentos e Inspiração Biológica}

Os Algoritmos Genéticos, sistematizados por \citeonline{Goldberg1989} e fundamentados no trabalho pioneiro de \citeonline{Holland1975}, são técnicas de busca e otimização inspiradas nos mecanismos da evolução biológica e na teoria da seleção natural de Charles Darwin. A premissa central é que, através de processos de seleção, reprodução e mutação, uma população de indivíduos evolui progressivamente em direção a soluções mais adaptadas ao ambiente.

Diferente de métodos determinísticos que iteram sobre um único ponto no espaço de busca, o AG opera sobre uma população de soluções candidatas simultaneamente. Cada indivíduo, denominado cromossomo, codifica uma solução potencial para o problema. No contexto da otimização de manobras, o genótipo do indivíduo é composto pelas variáveis de decisão (componentes do vetor impulso $\Delta \vec{v}$ e instante da manobra $t_m$), enquanto o fenótipo corresponde ao desempenho operacional dessa manobra.

A qualidade de cada solução é quantificada por uma função de aptidão (\textit{fitness function}), que avalia o quão bem o indivíduo atende aos objetivos da missão, como a minimização do consumo de propelente e a redução do risco de colisão. O processo evolutivo ocorre em ciclos geracionais: indivíduos mais aptos possuem maior probabilidade de sobreviver e transmitir seus genes, elevando a aptidão média da população ao longo das gerações.



\subsection{Codificação da Solução e Representação}

A etapa fundamental na aplicação de um AG é definir como a solução do problema será representada computacionalmente. Na formulação clássica de \citeonline{Holland1975}, os indivíduos são codificados como cadeias binárias (sequências de 0s e 1s), mimetizando a estrutura dos cromossomos biológicos. Nesta abordagem, uma variável contínua, como a velocidade de uma manobra, precisaria ser discretizada e convertida para binário.

No entanto, para problemas de engenharia de alta precisão como a astrodinâmica, a codificação binária apresenta desvantagens significativas:
\begin{enumerate}
    \item \textbf{Perda de Precisão:} A discretização impõe um limite na resolução da solução. Para obter a precisão necessária em uma manobra orbital (na ordem de milímetros por segundo), seriam necessárias cadeias binárias excessivamente longas.
    \item \textbf{Custo Computacional:} Exige processos constantes de codificação e decodificação a cada avaliação de função.
\end{enumerate}

Para superar essas limitações, este trabalho adota a Codificação Real (ou representação em ponto flutuante). Nesta abordagem, o ``cromossomo'' é um vetor de números reais que correspondem diretamente às grandezas físicas do problema, sem necessidade de tradução.

\begin{equation}
    \text{Indivíduo (Cromossomo)} = \mathbf{x} = [\Delta v_t, \Delta v_n, \Delta v_h, t_m]
\end{equation}

Onde cada gene representa o valor exato da componente de velocidade ou do tempo, respeitando a precisão de ponto flutuante da máquina (tipicamente \textit{double precision}). Isso permite que o algoritmo explore o espaço de busca contínuo com a sensibilidade necessária para o ajuste fino da trajetória.

\subsection{Operadores Genéticos para Codificação Real}

Com a adoção da representação em ponto flutuante, os operadores genéticos tradicionais de troca de bits tornam-se inadequados. Em seu lugar, aplicam-se operadores aritméticos desenvolvidos especificamente para espaços de busca contínuos, garantindo a integridade numérica da solução:

\begin{itemize}
    \item \textbf{Seleção:} Escolhe os indivíduos que atuarão como progenitores. O método de \textit{Torneio Binário} é utilizado, onde dois indivíduos são escolhidos aleatoriamente da população e o mais apto vence. Este método é preferido por permitir o ajuste da pressão seletiva sem exigir a ordenação completa da população.
    
    \item \textbf{Cruzamento (\textit{Crossover}):} Para a recombinação genética em domínio real, utiliza-se o operador SBX (\textit{Simulated Binary Crossover}), proposto por \citeonline{Deb1995}. Este operador foi desenhado para simular, em variáveis contínuas, as propriedades estatísticas de busca do crossover binário de um ponto. Ele gera descendentes cuja distribuição de probabilidade preserva a variância média dos pais, permitindo tanto a exploração próxima (filhos parecidos com os pais) quanto distante, sem a necessidade de discretização da variável.
    
    \item \textbf{Mutação:} Atua introduzindo pequenas perturbações estocásticas nos genes. Utiliza-se a \textit{Mutação Polinomial}, também formalizada no contexto de algoritmos reais por \citeonline{Deb1995} e \citeonline{Deb2001}. Ela adiciona um valor perturbativo governado por uma distribuição polinomial, essencial para evitar o estancamento em ótimos locais e garantir a ergodicidade da busca no espaço contínuo.
\end{itemize}

\subsection{Tratamento de Restrições em Otimização Evolutiva}

Em problemas de engenharia real, a otimização raramente é irrestrita. O objetivo de minimizar o custo ($\Delta v$) deve obedecer a requisitos estritos de segurança ($P_c \leq P_{lim}$). A adaptação de AGs para ambientes restritos é um desafio, uma vez que os operadores estocásticos podem gerar soluções inviáveis.

Uma abordagem robusta e consolidada na literatura recente é o método de tratamento de restrições baseado em regras de viabilidade, proposto por \citeonline{Deb2000}. Conhecida como \textit{Parameter-less Constraint Handling}, esta técnica dispensa o uso de pesos de penalidade (que requerem calibração sensível) ao modificar o critério de comparação no operador de seleção. Segundo as regras de Deb, a dominância entre duas soluções $x_1$ e $x_2$ é definida da seguinte forma:

\begin{enumerate}
    \item Se a solução $x_1$ é viável e $x_2$ é inviável, $x_1$ é preferida;
    \item Se ambas são inviáveis, prefere-se aquela com a menor violação das restrições;
    \item Se ambas são viáveis, prefere-se aquela com o melhor valor da função objetivo.
\end{enumerate}

Essa hierarquia garante que a busca pela eficiência energética ocorra prioritariamente dentro da região segura do espaço de busca, sendo ideal para aplicações críticas de segurança espacial.

\subsection{Aplicabilidade em Astrodinâmica}

A aplicação de Algoritmos Genéticos em problemas de astrodinâmica justifica-se pela complexidade das equações de movimento perturbado e pela natureza da função de probabilidade de colisão. Como o AG é um método de "caixa preta" e livre de derivadas (\textit{derivative-free}), ele permite o acoplamento direto com propagadores numéricos de alta fidelidade e modelos analíticos complexos (como o método de Chan), onde a relação explícita entre os parâmetros de manobra e o gradiente do risco não está disponível ou é computacionalmente custosa para ser obtida analiticamente \cite{Vasile2011}.

\section{Otimização de Manobras de Evasão de Colisão (CAMO)}
\label{sec:camo_review}

A literatura sobre Otimização de Manobras de Evasão de Colisão (CAMO) evoluiu de problemas simplificados de mudança orbital para sistemas complexos de decisão sob incerteza. Diferente das manobras de manutenção de órbita (\textit{station-keeping}), onde o objetivo é restaurar elementos orbitais nominais, a CAMO exige a satisfação de uma restrição de segurança probabilística ($P_c \le P_{lim}$) com o mínimo impacto na vida útil do satélite. A seguir, discute-se as principais vertentes de resolução deste problema.

\subsection{Tipologia da Manobra: Impulsiva versus Empuxo Contínuo}

A primeira grande divisão na literatura reside na modelagem da propulsão. A abordagem clássica considera manobras impulsivas (mudança instantânea de velocidade, $\Delta \vec{v}$), modelada idealmente para propulsores químicos de alto empuxo. \citeonline{Petropoulos2025} e \citeonline{Bombardelli2015} demonstram que, para avisos de colisão de curto prazo (24h a 48h antes do TCA), a manobra impulsiva é a estratégia mais eficiente energeticamente, pois permite explorar o efeito Oberth.

Conforme descrito por \citeonline{Sutton2016}, este efeito estabelece que a eficiência propulsiva na alteração da energia mecânica específica de uma órbita é proporcional à velocidade do veículo no momento da queima. Matematicamente, a variação da energia cinética é dada pela diferencial $dE = v dv$. Portanto, aplicar um $\Delta v$ quando o satélite possui alta velocidade escalar ($v$), como é o caso em LEO ($v \approx 7,5$ km/s), resulta em uma mudança de energia orbital significativamente maior do que a mesma aplicação em velocidades menores. Isso permite que pequenas correções impulsivas modifiquem o semi-eixo maior e, consequentemente, o período orbital, acumulando uma defasagem de posição (fase) suficiente para evitar a colisão após poucas revoluções.

Em contraste, com a popularização da propulsão elétrica, cresceu o interesse em manobras de baixo empuxo (\textit{low-thrust}). \citeonline{Peterson2020} e \citeonline{vittori2022} exploram o uso de arcos de empuxo contínuo. Embora estas soluções sejam extremamente eficientes em termos de consumo de massa — devido ao elevado Impulso Específico ($I_{sp}$), parâmetro que mensura a eficiência com que o motor converte o fluxo de propelente em força de empuxo \cite{Sutton2016} —, \citeonline{lee2017} alertam que a otimização de baixo empuxo introduz uma complexidade computacional severa associada à resolução do \textit{Two-Point Boundary Value Problem} (TPBVP).

Este problema consiste em determinar a trajetória dinâmica que conecta dois estados fixos (o estado inicial do satélite e a condição final de segurança) em um tempo finito. Matematicamente, o TPBVP não possui solução analítica fechada para dinâmicas perturbadas, exigindo métodos iterativos de ``disparo'' (\textit{shooting methods}) que são altamente sensíveis à estimativa inicial e computacionalmente onerosos. Além disso, perde-se a vantagem da mudança instantânea de geometria fornecida pelo efeito Oberth. Devido à urgência típica dos alertas de colisão, esta dissertação foca no paradigma impulsivo.

\subsection{Limitações das Abordagens Geométricas}

Historicamente, o problema foi tratado de forma geométrica: buscava-se maximizar a distância de separação (\textit{miss distance}) no plano de encontro. \citeonline{Patera2001} e \citeonline{patera2003} formalizaram métodos para calcular o menor $\Delta V$ necessário para atingir uma separação alvo.

Contudo, a maximização da distância geométrica baseia-se na premissa implícita de que a incerteza posicional é esférica ou isotrópica. \citeonline{Alfano2005} e \citeonline{Conway2001} refutaram a eficácia dessa abordagem para casos realistas, onde as matrizes de covariância são altamente alongadas na direção da velocidade (formato de ``charuto''). \citeonline{Bombardelli2015} provou analiticamente que uma manobra que maximiza a distância pode, contra-intuitivamente, aumentar a probabilidade de colisão se deslocar o objeto secundário para o eixo maior da elipse de incerteza do primário. Isso consolidou a necessidade de usar a Probabilidade de Colisão ($P_c$) diretamente como métrica de otimização, e não a distância.

\subsection{O Debate: Otimização Convexa vs. Não-Convexa}

Para lidar com a restrição de $P_c$, uma vertente da literatura buscou ``convexificar'' o problema. \citeonline{Armellin2021} e \citeonline{Benedikter2019} propuseram transformações que tornam o espaço de busca convexo, permitindo o uso de solvers de segunda ordem (SOCP) com garantia de convergência global em milissegundos. \citeonline{Dutta2022} expandiram essa técnica para incluir incertezas limitadas na execução.

Apesar da elegância matemática, a convexificação impõe linearizações na dinâmica relativa e simplificações no modelo de risco. Em uma análise crítica comparativa, \citeonline{dutta2023a} demonstrou que métodos convexos tendem a falhar em cenários onde a dinâmica perturbada (arrasto diferencial intenso) ou a não-gaussianidade das incertezas dominam. \citeonline{Gonzalo2021} reforça que, para órbitas com alta excentricidade ou encontros de baixa velocidade relativa, a topologia do risco é inerentemente não-convexa e multimodal, exigindo métodos globais.

\subsection{Algoritmos Evolutivos e Multiobjetividade}

Dada a complexidade descrita, meta-heurísticas consolidaram-se como a abordagem robusta por excelência. \citeonline{Lee2012} aplicaram Algoritmos Genéticos (AG) acoplados diretamente à fórmula de \citeonline{Chan1997}, demonstrando que o AG consegue navegar por superfícies de risco desconexas onde métodos de gradiente estagnam.

A grande vantagem dos métodos evolutivos, contudo, reside na capacidade multiobjetivo. Em operações reais, a ``melhor'' manobra não é apenas a de menor $\Delta V$, mas aquela que equilibra múltiplos fatores. \citeonline{Vasile2011} e \citeonline{Morselli2015} utilizaram MOEAs (\textit{Multi-Objective Evolutionary Algorithms}) para gerar Fronteiras de Pareto, revelando o \textit{trade-off} entre consumo e segurança. \citeonline{Meskini2014} adicionou uma camada extra de realismo ao co-otimizar a evasão com a manutenção da posição orbital (station-keeping), evitando que a manobra de desvio prejudique a missão principal do satélite.

Mais recentemente, \citeonline{Zhu2023} aplicaram algoritmos meméticos (híbridos) para refinar essas soluções, enquanto \citeonline{Nazari2022} explorou o algoritmo Firefly. Embora estas novas meta-heurísticas apresentem resultados promissores em benchmarks acadêmicos, a aplicação em operações reais de segurança favorece técnicas com comportamento de convergência amplamente documentado.

Neste contexto, a base teórica consolidada dos AGs \cite{Goldberg1989}, combinada com o método de tratamento de restrições de \citeonline{Deb2000}, confere a esta abordagem uma maturidade técnica superior. Diferente de algoritmos inspirados na natureza mais recentes, cujos hiperparâmetros de ajuste e estabilidade em topologias complexas ainda são objeto de investigação ativa, o AG oferece um compromisso comprovado entre exploração global e robustez, justificando sua escolha como a ferramenta mais adequada para o escopo crítico deste trabalho.

\subsection{Robustez sob Incerteza de Execução}

Por fim, a literatura atual move-se da otimização nominal para a robusta. Uma solução ótima matemática pode ser operacionalmente perigosa se for sensível a pequenos erros de execução. \citeonline{Wang2023} abordam o planejamento considerando erros de magnitude e apontamento do empuxo (\textit{thrust errors}). \citeonline{Sala2018} utiliza expansão polinomial do caos para propagar incertezas de navegação dentro do loop de otimização. \citeonline{dutta2023b} investiga como a natureza da incerteza (epistêmica vs. aleatória) altera a topologia da solução ótima.

\section{Síntese e Perspectivas}
\label{sec:sintese}

Este capítulo apresentou uma revisão abrangente da fundamentação teórica necessária para abordar o problema de otimização de manobras de evasão de colisão. Partindo dos fundamentos da mecânica orbital, passando pela análise de risco e chegando às estratégias de otimização, procurou-se estabelecer uma base sólida que integra conceitos de múltiplas disciplinas: dinâmica orbital, modelagem estocástica e otimização matemática.

A literatura sobre CAMO é vasta e multifacetada, refletindo a importância crescente do problema. Observa-se uma clara divisão entre abordagens analíticas e convexas, que oferecem elegância matemática e eficiência computacional, e abordagens heurísticas, que oferecem flexibilidade e robustez para lidar com problemas complexos e não-convexos. Cada abordagem possui seus méritos e limitações, e a escolha apropriada depende fortemente do contexto operacional e das restrições computacionais disponíveis.

A presente dissertação se insere neste panorama ao propor uma análise sistemática com complexidade crescente, desde modelos simplificados até modelos de alta fidelidade com propagação numérica e restituição de órbita. Esta progressão permite não apenas explorar diferentes estratégias de otimização, mas também quantificar os trade-offs envolvidos em cada nível de fidelidade e complexidade. Os detalhes da metodologia adotada, bem como a justificativa das escolhas específicas de algoritmos e modelos, serão apresentados no capítulo seguinte.
